\documentclass[11pt]{article}
\usepackage[margin=18mm,headheight=22pt]{geometry}
\usepackage{amsmath}
\usepackage{fontspec}
\usepackage{unicode-math}
\setmainfont{NewCM10}[Extension=.otf,UprightFont=*-Book,BoldFont=*-Bold,ItalicFont=*-BookItalic,BoldItalicFont=*-BoldItalic]
\setsansfont{NewCMSans10}[Extension=.otf,UprightFont=*-Book,BoldFont=*-Bold,ItalicFont=*-BookOblique,BoldItalicFont=*-BoldOblique]
\setmathfont{NewCMMath-Book.otf}
\usepackage{xcolor}
\definecolor{FAUblue}{HTML}{04316A}
\definecolor{Paperblue}{HTML}{EEF2F7}
\definecolor{EEIgold}{HTML}{CC9C42}
\usepackage{etoolbox,enumitem,booktabs,tabularx,array,needspace,microtype}
\usepackage[most]{tcolorbox}
\usepackage{titlesec,fancyhdr}
\usepackage[hidelinks]{hyperref}
\usepackage{draftwatermark}
\SetWatermarkText{Made by Thanate Amornjiraporn (Focus)}
\SetWatermarkColor{black!7}\SetWatermarkFontSize{30pt}\SetWatermarkAngle{45}
\setlength{\parindent}{0pt}
\setlength{\parskip}{5pt}
\setlist{itemsep=2pt,topsep=3pt,leftmargin=1.6em,beginpenalty=10000,endpenalty=10000}
\raggedbottom
\clubpenalty=10000 \widowpenalty=10000 \displaywidowpenalty=10000 \interlinepenalty=10000
\setcounter{secnumdepth}{0}
\setcounter{tocdepth}{1}
\pagestyle{fancy}\fancyhf{}
\fancyhead[L]{\small\sffamily\color{FAUblue} From Fundamentals to AI Apocalypse · APDR · Blackboard Exercise 2}
\fancyhead[R]{\small\sffamily Blackboard Exercise 2: Worked solutions}
\fancyfoot[C]{\small\thepage}
\titleformat{\section}{\Large\bfseries\sffamily\color{FAUblue}}{}{0pt}{}
\titleformat{\subsection}{\large\bfseries\sffamily\color{FAUblue}}{}{0pt}{}
\titlespacing*{\section}{0pt}{16pt}{6pt}
\preto\section{\Needspace{8\baselineskip}}
\preto\subsection{\Needspace{6\baselineskip}}
\newtcolorbox{equationcard}[1]{colback=Paperblue,colframe=Paperblue,boxrule=0pt,arc=0pt,left=10pt,right=10pt,top=6pt,bottom=6pt,title={#1},fonttitle=\small\bfseries\sffamily,coltitle=FAUblue,before skip=8pt,after skip=8pt}
\newtcolorbox{workedbox}[1]{colback=white,colframe=EEIgold,boxrule=0pt,leftrule=3pt,arc=0pt,left=10pt,top=4pt,bottom=4pt,title={#1},fonttitle=\small\bfseries\sffamily,coltitle=FAUblue,attach title to upper={\par},before skip=8pt,after skip=8pt}
\newtcolorbox{recallbox}[1]{colback=white,colframe=FAUblue,boxrule=0pt,leftrule=3pt,arc=0pt,left=10pt,top=4pt,bottom=4pt,title={#1},fonttitle=\small\bfseries\sffamily,coltitle=FAUblue,attach title to upper={\par},before skip=8pt,after skip=8pt}
\usepackage{environ}
\newsavebox\measurebox
% Typeset the section once invisibly to measure it; start a new page if it will not fit in what is left.
\makeatletter
\newif\ifkeep@measuring
\NewEnviron{keepsection}{%
  \ifkeep@measuring \BODY \else
    \edef\keep@sec{\the\value{section}}\edef\keep@sub{\the\value{subsection}}%
    \setbox\measurebox\vbox{\hsize\linewidth\keep@measuringtrue\def\Needspace##1{}\BODY}%
    \setcounter{section}{\keep@sec}\setcounter{subsection}{\keep@sub}%
    \dimen2=\ht\measurebox \advance\dimen2\dp\measurebox
    \dimen@=\textheight
    \ifdim\pagegoal<\maxdimen
      \dimen@=\pagegoal \advance\dimen@-\pagetotal
      \ifdim\dimen2>\dimen@ \clearpage \dimen@=\textheight \fi
    \fi
    % if it fits, nothing inside may force an early break
    \ifdim\dimen2>\dimen@ \BODY \else \begingroup\def\Needspace##1{}\BODY\endgroup \fi
  \fi}
\makeatother
\newtcolorbox{notebox}{colback=Paperblue,colframe=Paperblue,boxrule=0pt,arc=0pt,left=8pt,right=8pt,top=4pt,bottom=4pt,before skip=6pt,after skip=6pt}
\newtcolorbox{plainbox}{colback=white,colframe=FAUblue,boxrule=0pt,leftrule=3pt,arc=0pt,left=10pt,right=8pt,top=4pt,bottom=4pt,before skip=6pt,after skip=8pt}
\usetikzlibrary{arrows.meta}
\title{\sffamily\bfseries\color{FAUblue}Blackboard Exercise 2: Worked solutions}
\author{Algorithms, Programming, and Data Representation \ensuremath{\cdot} Module 92313 \ensuremath{\cdot} Blackboard Exercise 2\\[2pt]\small From Fundamentals to AI Apocalypse · personal study guide}
\date{}
\begin{document}
\maketitle
\begingroup\footnotesize\setlength{\parskip}{0pt}\tableofcontents\endgroup
\clearpage

\Needspace{8\baselineskip}
{\large A worked route through number bases, binary fractions, integer encodings, ASCII and Boolean algebra. Try each subpart before opening its steps.}


\begin{keepsection}
\Needspace{5\baselineskip}
\section{Exercise 2.1: Natural numbers}

\textbf{Connected topic:} Number systems and data representation (§3.1–3.3)

\begin{keepsection}
\Needspace{7\baselineskip}
\subsection{2.1.1 Number representations}

\Needspace{4\baselineskip}
Try first: Complete the five conversions; show how the first and last binary rows convert directly to octal and hexadecimal.

\nopagebreak[4]{\small\itshape Answers: see “Exercise 2.1: Natural numbers” under Answers at the end.}
\end{keepsection}

\begin{keepsection}
\Needspace{7\baselineskip}
\subsection{2.1.2 Range of an unsigned representation}

\Needspace{4\baselineskip}
Try first: Find the ranges for 8, 16, and 32 bits, derive the general rule, then find the fewest bits that cover 0 through 1000.

\nopagebreak[4]{\small\itshape Answers: see “Exercise 2.1: Natural numbers” under Answers at the end.}
\end{keepsection}

\begin{keepsection}
\Needspace{7\baselineskip}
\subsection{2.1.3 Representation and storage}

\Needspace{4\baselineskip}
Try first: Decide which listed values fit an unsigned byte, identify equal values, and explain why distinct spellings can denote the same number.

\nopagebreak[4]{\small\itshape Answers: see “Exercise 2.1: Natural numbers” under Answers at the end.}
\end{keepsection}

\begin{keepsection}
\Needspace{7\baselineskip}
\subsection{2.1.4 Arithmetic}

\Needspace{4\baselineskip}
Try first: Calculate the four sums, difference, and products in the operand bases without converting the operands to decimal first; decide whether each result fits an unsigned byte.

\nopagebreak[4]{\small\itshape Answers: see “Exercise 2.1: Natural numbers” under Answers at the end.}
\end{keepsection}

\begin{keepsection}
\Needspace{7\baselineskip}
\subsection{2.1.5 Overflow}

\Needspace{4\baselineskip}
Try first: For each decimal addition, find the mathematical sum, whether an 8-bit unsigned overflow occurs, and the stored byte. Then explain the NumPy result for uint8(250)+uint8(13).

\nopagebreak[4]{\small\itshape Answers: see “Exercise 2.1: Natural numbers” under Answers at the end.}
\end{keepsection}

\begin{keepsection}
\Needspace{7\baselineskip}
\subsection{2.1.6 Working with an unknown radix}

\Needspace{4\baselineskip}
Try first: The numeral 243 in base r has decimal value 99. Find r and decide whether 354\textsubscript{r} and 358\textsubscript{r} are valid.

\nopagebreak[4]{\small\itshape Answers: see “Exercise 2.1: Natural numbers” under Answers at the end.}
\end{keepsection}
\end{keepsection}

\begin{keepsection}
\Needspace{5\baselineskip}
\section{Exercise 2.2: Floating-point numbers}

\textbf{Connected topic:} Binary fractions and floating-point representation (§4.1)

\begin{keepsection}
\Needspace{7\baselineskip}
\subsection{2.2.1 Binary fractions}

\Needspace{4\baselineskip}
Try first: Complete the four binary/decimal conversions and show one calculation in each direction.

\nopagebreak[4]{\small\itshape Answers: see “Exercise 2.2: Floating-point numbers” under Answers at the end.}
\end{keepsection}

\begin{keepsection}
\Needspace{7\baselineskip}
\subsection{2.2.2 Finite and repeating representations}

\Needspace{4\baselineskip}
Try first: Convert 0.1\textsubscript{10} until its repeating pattern is visible. Decide which listed fractions terminate in binary, then state the general condition.

\nopagebreak[4]{\small\itshape Answers: see “Exercise 2.2: Floating-point numbers” under Answers at the end.}
\end{keepsection}
\end{keepsection}

\begin{keepsection}
\Needspace{5\baselineskip}
\section{Exercise 2.3: Signed and unsigned integers}

\textbf{Connected topic:} Integer representations and arithmetic (§3.4)

\begin{keepsection}
\Needspace{7\baselineskip}
\subsection{2.3.1 Two’s complement}

\Needspace{4\baselineskip}
Try first: Complete the table for the three bit patterns and the two decimal values, using exactly eight bits.

\nopagebreak[4]{\small\itshape Answers: see “Exercise 2.3: Signed and unsigned integers” under Answers at the end.}
\end{keepsection}

\begin{keepsection}
\Needspace{7\baselineskip}
\subsection{2.3.2 Range and interpretation}

\Needspace{4\baselineskip}
Try first: Find the 8-bit and general n-bit ranges, explain the asymmetric counts, and interpret 11110110\textsubscript{2} both ways.

\nopagebreak[4]{\small\itshape Answers: see “Exercise 2.3: Signed and unsigned integers” under Answers at the end.}
\end{keepsection}

\begin{keepsection}
\Needspace{7\baselineskip}
\subsection{2.3.3 Signed overflow}

\Needspace{4\baselineskip}
Try first: Add 01100101\textsubscript{2} and 00111110\textsubscript{2}, interpret all three patterns as signed decimal, and test for signed overflow.

\nopagebreak[4]{\small\itshape Answers: see “Exercise 2.3: Signed and unsigned integers” under Answers at the end.}
\end{keepsection}
\end{keepsection}

\begin{keepsection}
\Needspace{5\baselineskip}
\section{Exercise 2.4: Text representation and interpretation}

\textbf{Connected topic:} Text and character representation (§3.5)

\begin{keepsection}
\Needspace{6\baselineskip}
\subsection{2.4.1 ASCII}

\Needspace{3\baselineskip}
Try first: Complete the table and encode AIBE as four 8-bit ASCII bytes.

\nopagebreak[4]{\small\itshape Answers: see “Exercise 2.4: Text representation and interpretation” under Answers at the end.}
\end{keepsection}

\begin{keepsection}
\Needspace{7\baselineskip}
\subsection{2.4.2 One bit pattern, different meanings}

\Needspace{4\baselineskip}
Try first: Interpret 01000001\textsubscript{2} as unsigned, as 8-bit two’s complement, and as ASCII. State what extra information a reader needs.

\nopagebreak[4]{\small\itshape Answers: see “Exercise 2.4: Text representation and interpretation” under Answers at the end.}
\end{keepsection}
\end{keepsection}

\begin{keepsection}
\Needspace{5\baselineskip}
\section{Exercise 2.5: Boolean algebra}

\textbf{Connected topic:} Boolean algebra identities (§1.1–1.2)

\begin{keepsection}
\Needspace{6\baselineskip}
\subsection{2.5.1 Truth table}

\Needspace{3\baselineskip}
Try first: Fill the columns A\ensuremath{{}^{\prime}}, A\ensuremath{\cdot}B, A+B, and A+(B\ensuremath{\cdot}C) for all input rows in the given order.

\nopagebreak[4]{\small\itshape Answers: see “Exercise 2.5: Boolean algebra” under Answers at the end.}
\end{keepsection}

\begin{keepsection}
\Needspace{6\baselineskip}
\subsection{2.5.2 Simplification}

\Needspace{3\baselineskip}
Try first: Simplify the four expressions and name the Boolean law used at each step.

\nopagebreak[4]{\small\itshape Answers: see “Exercise 2.5: Boolean algebra” under Answers at the end.}
\end{keepsection}

\begin{keepsection}
\Needspace{7\baselineskip}
\subsection{2.5.3 Equivalence and construction}

\Needspace{4\baselineskip}
Try first: Compare A+(B\ensuremath{\cdot}C) with (A+B)(A+C) using algebra and a truth table; construct F=1 exactly when at least two inputs are 1.

\nopagebreak[4]{\small\itshape Answers: see “Exercise 2.5: Boolean algebra” under Answers at the end.}
\end{keepsection}
\end{keepsection}


\clearpage
\section{Answers}
\begin{keepsection}\subsection{Exercise 2.1: Natural numbers}

\begin{workedbox}{2.1.1 Number representations \ensuremath{\cdot} worked solution}
\Needspace{10\baselineskip}
\begin{enumerate}
\item Read each subscript as the base. For binary-to-octal, group bits from the binary point in threes; for binary-to-hexadecimal, group in fours. Pad only on the left with zeroes when needed.
\item 10110110\textsubscript{2} = 010 110 110\textsubscript{2} = 266\textsubscript{8}, and 1011 0110\textsubscript{2} = B6\textsubscript{16}. The complete table is: 10110110\textsubscript{2} = 266\textsubscript{8} = 182\textsubscript{10} = B6\textsubscript{16}; 347\textsubscript{8} = 11100111\textsubscript{2} = 231\textsubscript{10} = E7\textsubscript{16}; 156\textsubscript{10} = 10011100\textsubscript{2} = 234\textsubscript{8} = 9C\textsubscript{16}; 2D7\textsubscript{16} = 0010 1101 0111\textsubscript{2} = 1327\textsubscript{8} = 727\textsubscript{10}; 110101011001\textsubscript{2} = 6531\textsubscript{8} = D59\textsubscript{16} = 3417\textsubscript{10}.
\item The groups work because 8 = 2\textsuperscript{3} and 16 = 2\textsuperscript{4}: each octal digit encodes exactly three bits, and each hexadecimal digit exactly four. Check by expanding any one row back into powers of its base.
\end{enumerate}
\end{workedbox}

\end{keepsection}
\begin{keepsection}\subsection{Exercise 2.1: Natural numbers}

\begin{workedbox}{2.1.2 Range of an unsigned representation \ensuremath{\cdot} worked solution}
\Needspace{7\baselineskip}
\begin{enumerate}
\item An n-bit unsigned pattern has place values 2\textsuperscript{0} through 2\textsuperscript{n\ensuremath{-}1}. Its minimum is all zeroes; its maximum is their sum.
\item The ranges are 0–255, 0–65,535, and 0–4,294,967,295. In general: 0 \ensuremath{\leq} x \ensuremath{\leq} 2\textsuperscript{n}\ensuremath{-}1. For 0 through 1000, 2\textsuperscript{9}\ensuremath{-}1=511 is too small, while 2\textsuperscript{10}\ensuremath{-}1=1023 is enough, so 10 bits are required.
\item The count of distinct patterns is 2\textsuperscript{n}. Verify minimality by showing 2\textsuperscript{9} \ensuremath{<} 1001 \ensuremath{\leq} 2\textsuperscript{10}.
\end{enumerate}
\end{workedbox}

\end{keepsection}
\begin{keepsection}\subsection{Exercise 2.1: Natural numbers}

\begin{workedbox}{2.1.3 Representation and storage \ensuremath{\cdot} worked solution}
\Needspace{9\baselineskip}
\begin{enumerate}
\item An unsigned byte stores 0 through 255. Convert each representation to its value, then compare with that range.
\item 11111111\textsubscript{2}=255, 10000000\textsubscript{2}=128, FF\textsubscript{16}=255, and 256\textsubscript{10}=256. The first and third denote the same value and both fit; 128 fits; 256 does not.
\item A numeral's value is the weighted sum of its digits in its base. Base-2 11111111 and base-16 FF are different spellings but both sum to 255. A byte stores the value's bit pattern, not the textual spelling.
\end{enumerate}
\end{workedbox}

\end{keepsection}
\begin{keepsection}\subsection{Exercise 2.1: Natural numbers}

\begin{workedbox}{2.1.4 Arithmetic \ensuremath{\cdot} worked solution}
\Needspace{9\baselineskip}
\begin{enumerate}
\item Keep the radix fixed during each operation, align place values, and perform carries or borrows in that radix.
\item 10110110\textsubscript{2} + 00101101\textsubscript{2} = 11100011\textsubscript{2} (227), which fits. 11001010\textsubscript{2} \ensuremath{-} 00110111\textsubscript{2} = 10010011\textsubscript{2} (147), which fits. 00101101\textsubscript{2} \ensuremath{\times} 00000100\textsubscript{2} = 10110100\textsubscript{2} (180), which fits. A7\textsubscript{16} + 3C\textsubscript{16} = E3\textsubscript{16} (227), which fits.
\item All four values are at most FF\textsubscript{16}=255. The subtraction is nonnegative; the multiplication by 4 is a two-bit left shift, which provides a quick check.
\end{enumerate}
\end{workedbox}

\end{keepsection}
\begin{keepsection}\subsection{Exercise 2.1: Natural numbers}

\begin{workedbox}{2.1.5 Overflow \ensuremath{\cdot} worked solution}
\Needspace{10\baselineskip}
\begin{enumerate}
\item An 8-bit unsigned value stores the sum only when it is at most 255. For modular storage, retain the low eight bits, equivalently take the remainder modulo 256.
\item 200+55=255: no overflow, stored 255. 200+56=256: overflow, stored 0. 250+13=263: overflow, stored 7. 128+128=256: overflow, stored 0. The condition is a+b\ensuremath{>}255. For the NumPy expression, the expected 8-bit result is 7 because 263 mod 256=7; this is fixed-width unsigned arithmetic; the NumPy 2.5.4 run also emits a RuntimeWarning for scalar overflow.
\item Compare each stored value with its mathematical sum modulo 256. The boundary case 200+55 is exactly 255, so it does not overflow.
\end{enumerate}
\end{workedbox}

\end{keepsection}
\begin{keepsection}\subsection{Exercise 2.1: Natural numbers}

\begin{workedbox}{2.1.6 Working with an unknown radix \ensuremath{\cdot} worked solution}
\Needspace{6\baselineskip}
\begin{enumerate}
\item Expand by place value: 243\textsubscript{r} = 2r\textsuperscript{2}+4r+3. A radix must be an integer larger than every digit.
\item Solve 2r\textsuperscript{2}+4r+3=99, giving r\textsuperscript{2}+2r\ensuremath{-}48=0=(r+8)(r\ensuremath{-}6). The positive radix is r=6. In base six, 354\textsubscript{6} is valid because each digit is below 6; 358\textsubscript{6} is invalid because digit 8 is not allowed.
\item Substitute r=6: 2\ensuremath{\cdot}36+4\ensuremath{\cdot}6+3=99. Reject \ensuremath{-}8 because a radix is positive.
\end{enumerate}
\end{workedbox}

\end{keepsection}
\begin{keepsection}\subsection{Exercise 2.2: Floating-point numbers}

\begin{workedbox}{2.2.1 Binary fractions \ensuremath{\cdot} worked solution}
\Needspace{8\baselineskip}
\begin{enumerate}
\item After the binary point, places have weights 2\textsuperscript{\ensuremath{-}}\textsuperscript{1}, 2\textsuperscript{\ensuremath{-}}\textsuperscript{2}, and so on. To convert a binary fraction, add the place values of its 1 bits.
\item 0.1011\textsubscript{2}=1/2+1/8+1/16=0.6875\textsubscript{10}. 10.011\textsubscript{2}=2+1/4+1/8=2.375\textsubscript{10}. For the reverse direction, repeatedly double the fractional part: 0.625\ensuremath{\times}2=1.25 gives bit 1 and remainder .25; .25\ensuremath{\times}2=.5 gives 0; .5\ensuremath{\times}2=1 gives 1. Thus 0.625\textsubscript{10}=0.101\textsubscript{2}. Similarly 5.375\textsubscript{10}=101.011\textsubscript{2}.
\item Convert back: 0.101\textsubscript{2}=1/2+1/8=0.625, and 101.011\textsubscript{2}=4+1+1/4+1/8=5.375.
\end{enumerate}
\end{workedbox}

\end{keepsection}
\begin{keepsection}\subsection{Exercise 2.2: Floating-point numbers}

\begin{workedbox}{2.2.2 Finite and repeating representations \ensuremath{\cdot} worked solution}
\Needspace{9\baselineskip}
\begin{enumerate}
\item Repeatedly multiply the fractional remainder by 2; each product's integer part is the next bit.
\item For 1/10: remainders produce 0.0001100110011…\textsubscript{2} = 0.0(0011)\textsubscript{2}, with the block 0011 repeating. Reduce the candidates: 1/2, 3/8, 1/10, 5/16, 7/20. The finite ones are 1/2, 3/8, and 5/16; 1/10 and 7/20 repeat.
\item A reduced fraction p/q has a finite binary expansion exactly when q is a power of 2. Here 10=2\ensuremath{\cdot}5 and 20=4\ensuremath{\cdot}5 retain a factor 5 after reduction, so those expansions cannot terminate.
\end{enumerate}
\end{workedbox}

\end{keepsection}
\begin{keepsection}\subsection{Exercise 2.3: Signed and unsigned integers}

\begin{workedbox}{2.3.1 Two’s complement \ensuremath{\cdot} worked solution}
\Needspace{8\baselineskip}
\begin{enumerate}
\item For a nonnegative value, write its ordinary binary form padded to eight bits. For a negative value, invert the bits of its magnitude and add one.
\item 01111111\textsubscript{2}=127; 11111111\textsubscript{2}=\ensuremath{-}1; 10000000\textsubscript{2}=\ensuremath{-}128. Conversely, \ensuremath{-}8 is 11111000\textsubscript{2}, and \ensuremath{-}37 is 11011011\textsubscript{2}.
\item Check negative encodings by interpreting the most significant bit as weight \ensuremath{-}128 and the remaining bits as positive powers: 11111000\textsubscript{2}=\ensuremath{-}128+120=\ensuremath{-}8; 11011011\textsubscript{2}=\ensuremath{-}128+64+16+8+2+1=\ensuremath{-}37.
\end{enumerate}
\end{workedbox}

\end{keepsection}
\begin{keepsection}\subsection{Exercise 2.3: Signed and unsigned integers}

\begin{workedbox}{2.3.2 Range and interpretation \ensuremath{\cdot} worked solution}
\Needspace{9\baselineskip}
\begin{enumerate}
\item In n-bit two’s complement, the leading bit has weight \ensuremath{-}2\textsuperscript{n\ensuremath{-}1}; the other bits sum from 0 to 2\textsuperscript{n\ensuremath{-}1}\ensuremath{-}1.
\item The 8-bit range is \ensuremath{-}128 through 127. Generally it is \ensuremath{-}2\textsuperscript{n\ensuremath{-}1} through 2\textsuperscript{n\ensuremath{-}1}\ensuremath{-}1. There are 2\textsuperscript{n\ensuremath{-}1} negative patterns but only 2\textsuperscript{n\ensuremath{-}1}\ensuremath{-}1 positive patterns because zero uses one of the nonnegative patterns. 11110110\textsubscript{2} is 246 unsigned and \ensuremath{-}10 signed (\ensuremath{-}128+64+32+16+4+2).
\item The same stored bits have no inherent signedness. The surrounding type or format decides whether the leading bit has weight +128 or \ensuremath{-}128.
\end{enumerate}
\end{workedbox}

\end{keepsection}
\begin{keepsection}\subsection{Exercise 2.3: Signed and unsigned integers}

\begin{workedbox}{2.3.3 Signed overflow \ensuremath{\cdot} worked solution}
\Needspace{8\baselineskip}
\begin{enumerate}
\item Add the bit patterns: 01100101\textsubscript{2}+00111110\textsubscript{2}=10100011\textsubscript{2}.
\item The operands are +101 and +62; the 8-bit result pattern 10100011\textsubscript{2} represents \ensuremath{-}93 (\ensuremath{-}128+32+2+1). The mathematical sum is 163, outside the range \ensuremath{-}128…127, so signed overflow occurs.
\item Two positive operands produced a negative result, which detects overflow. In general, signed addition overflows when operand signs match and the result sign differs; opposite-sign operands cannot overflow when added.
\end{enumerate}
\end{workedbox}

\end{keepsection}
\begin{keepsection}\subsection{Exercise 2.4: Text representation and interpretation}

\begin{workedbox}{2.4.1 ASCII \ensuremath{\cdot} worked solution}
\Needspace{8\baselineskip}
\begin{enumerate}
\item Treat each ASCII code as a byte. Convert binary to decimal and look up the corresponding standard ASCII character; reverse the process for a given decimal code.
\item 01000001\textsubscript{2}=65\textsubscript{10}='A'; 01100001\textsubscript{2}=97\textsubscript{10}='a'; 48\textsubscript{10}=00110000\textsubscript{2}='0'. AIBE is 01000001 01001001 01000010 01000101\textsubscript{2}, or 41 49 42 45 in hexadecimal.
\item Check that uppercase and lowercase letters differ by 32 in ASCII (one bit at this position), and that decimal 48 is the character '0', not the numeric value zero.
\end{enumerate}
\end{workedbox}

\end{keepsection}
\begin{keepsection}\subsection{Exercise 2.4: Text representation and interpretation}

\begin{workedbox}{2.4.2 One bit pattern, different meanings \ensuremath{\cdot} worked solution}
\Needspace{7\baselineskip}
\begin{enumerate}
\item Evaluate the same eight bits under each encoding rule rather than changing the bits.
\item Unsigned interpretation: 65. Two’s-complement interpretation: 65, because the leading bit is zero. ASCII interpretation: 'A'.
\item The data’s type or agreed encoding supplies the meaning: for example, unsigned byte, signed 8-bit integer, or ASCII character. Bits alone do not specify their interpretation.
\end{enumerate}
\end{workedbox}

\end{keepsection}
\begin{keepsection}\subsection{Exercise 2.5: Boolean algebra}

\begin{workedbox}{2.5.1 Truth table \ensuremath{\cdot} worked solution}
\Needspace{8\baselineskip}
\begin{enumerate}
\item Evaluate each row independently with 0=false and 1=true. Complement flips a bit; \ensuremath{\cdot} is AND; + is OR. Parentheses make B\ensuremath{\cdot}C happen before the outer OR.
\item In the row order 000,001,010,011,100,101,110,111, the columns are: A\ensuremath{{}^{\prime}} = 1,1,1,1,0,0,0,0; A\ensuremath{\cdot}B = 0,0,0,0,0,0,1,1; A+B = 0,0,1,1,1,1,1,1; A+(B\ensuremath{\cdot}C) = 0,0,0,1,1,1,1,1.
\item Spot-check 011: A=0 and B\ensuremath{\cdot}C=1, so the final expression is 1. For all rows with A=1, both OR columns are 1.
\end{enumerate}
\end{workedbox}

\end{keepsection}
\begin{keepsection}\subsection{Exercise 2.5: Boolean algebra}

\begin{workedbox}{2.5.2 Simplification \ensuremath{\cdot} worked solution}
\Needspace{7\baselineskip}
\begin{enumerate}
\item Use complement, identity, distributive, and absorption laws; show the intermediate rewrite so each result can be checked.
\item A+A\ensuremath{{}^{\prime}}=1 by complement. A\ensuremath{\cdot}A\ensuremath{{}^{\prime}}=0 by complement. A+(A\ensuremath{{}^{\prime}}\ensuremath{\cdot}B)=(A+A\ensuremath{{}^{\prime}})(A+B)=1\ensuremath{\cdot}(A+B)=A+B, by distributivity and identity. (A+B)(A+B\ensuremath{{}^{\prime}})=A+BB\ensuremath{{}^{\prime}}=A+0=A, by distributivity and complement.
\item Check the last two with A=0 and A=1 separately; their outputs match A+B and A respectively.
\end{enumerate}
\end{workedbox}

\end{keepsection}
\begin{keepsection}\subsection{Exercise 2.5: Boolean algebra}

\begin{workedbox}{2.5.3 Equivalence and construction \ensuremath{\cdot} worked solution}
\Needspace{9\baselineskip}
\begin{enumerate}
\item Expand the product using distributivity, then use idempotence and absorption. For the construction, list the input rows with at least two 1s.
\item (A+B)(A+C)=A\ensuremath{\cdot}A+A\ensuremath{\cdot}C+A\ensuremath{\cdot}B+B\ensuremath{\cdot}C=A+A\ensuremath{\cdot}C+A\ensuremath{\cdot}B+B\ensuremath{\cdot}C=A+B\ensuremath{\cdot}C. Thus the expressions are equivalent. For row order 000 through 111, both evaluate to 0,0,0,1,1,1,1,1.
\item The true rows for the “at least two” function are 011, 101, 110, and 111, giving F=A\ensuremath{{}^{\prime}}BC+AB\ensuremath{{}^{\prime}}C+ABC\ensuremath{{}^{\prime}}+ABC. Grouping pairs or using absorption yields F=AB+AC+BC. Check: exactly two 1s makes one pair term true; three 1s makes all three true.
\end{enumerate}
\end{workedbox}

\end{keepsection}

\vfill{\small Personal worked solutions based on the APDR Blackboard Exercise 2 sheet, written in my own words; no exercise-sheet prose or figures are reproduced. Not an official FAU publication; check course details against the current materials. Made by \textbf{Thanate Amornjiraporn (Focus)}.}

\end{document}
